\input{preamble}
\author{Черносов Вячеслав}
\title{Моделирование атмосферы нейтронной звезды с учетом слоя растекания}

\begin{document}
\maketitle
В моделировании атмосферы нейтронной звезды (НЗ) можно выделить три основных этапа: 
\begin{itemize}
    \item вычисление параметров НЗ
    \item вычисление параметров поверхности НЗ 
    \item вычисление параметров спектра НЗ
\end{itemize}
Этим мы и займемся в соответсвующем порядке.

\section{Нейтронная звезда}

Для начала, зададим основные параметры нейтронной звезды:
\begin{itemize}
    \item радиус звезды $R_{\star}$ 
    \item масса звезды $M_{\star}$
    \item частота вращения звезды $\nu_{\star}$ 
    \item угол наклона оси вращения звезды к наблюдателю $i$
    \item химический состав атмосферы звезды $C$
\end{itemize}
Эти параметры мы фиксируем в качестве постоянных для наших вычислений.

\subsection{Физические параметры}
Для начала разберемся с частотой вращения нейтронной звезды. Циклическая частота вращения может быть найдена как:
$$\Omega_{\star} = 2 \pi \nu_{\star}$$

Однако, существует физическое ограничение, связанное с уравнением состояния нейтронной звезды, которое дает нам критическую частоту вращения:
$$\nu_{cr} = 1278 \text{ Hz } \cdot \sqrt{\frac{M_{1.4}}{(R_{10})^{3}}},$$
где $R_{10} = R_{\star}/(10\text{ km})$, $M_{1.4} = M_{\star}/(1.4 M_{\odot})$, благодаря чему мы можем отнормировать частоту вращения на критическую:
$$\bar{\nu} = \frac{\nu_{\star}}{\nu_{cr}}.$$

C помощью этой отнормированной частоты вращения, а также уравнения состояния, мы можем получить значения скорректированной массы нейтронной звезды:
$$ M' =  M_{\star} \cdot \left(a_0 - \frac{a_1}{\bar{\nu} - 1.1} + a_2 \cdot \bar{\nu}^2 \right) $$
где $a_1 = 0.001 \cdot (M_{1.4})^{1.5} $,  $a_2 = 10 a_1$, $a_0 = 1 - a_1/1.1$, а также её экваториального радиуса:


$$ R_{eq} = R_{\star} \cdot \left(r_0 - \frac{r_1}{\bar{\nu} - 1.07} + r_2 \cdot \bar{\nu}^2 \right),$$
где $r_1 = 0.025$, $r_2 = 0.07 \cdot (M_{1.4})^{1.5}$, $r_0 = 0.9766$.

Вычислим параметры, связанные с гравитацией. Радиус Шврацшильда будет равен:
$$R_{s} = \frac{2 GM_{\star}}{c^2},$$
а соответсвующее ему гравитационное красное смещение будет равно:
$$1 + z = \frac{1}{\sqrt{1 - R_s / R_{\star}}}.$$
Среднее значение гравитации на поверхности будет равно:
$$g_{\star} = \frac{GM_{\star}}{R_{\star}^2} \cdot (1 + z).$$
\subsection{Фотометрические параметры}
Прозрачность при когерентном томпсоновском рассеянии электронов будет равна:
$$\kappa = 0.2 \cdot (1 + X_C),$$
где $X_C$ -- массовая доля водорада для соответсвующего химического состава C. 
Тогда Эддингтоновский поток будет равен:
$$F^{(Edd)}_{\star} = \frac{g_{\star}c}{\kappa},$$
а соответсвующая температура по закону Стефана-Больцмана будет равна:
$$T^{(Edd)}_{\star} = \frac{(F^{(Edd)}_{\star}/\sigma_{SB})^{1/4}}{1 + z},$$
или же в энергитических единицах:
$$\Theta^{(Edd)}_{\star} = k_B T^{(Edd)}_{\star}.$$
Светимость же в этом случае будет равна:
$$L^{(Edd)}_{\star} = 4 \pi R_{\star}^2 \cdot F^{(Edd)}_{\star}$$
\subsection{Метрические коэффициенты}
Введем два безразмерных метрических коэффициента:
$$\chi = \frac{GM'}{R_{eq}c^2}, $$
$$\bar{\Omega}_{\star} =  \Omega_{\star} \cdot \sqrt{\frac{R_{eq}^3}{G M'}}.$$
C помощью них вычислим обезразмеренный момент инерции:
$$\bar{I} = \sqrt{\chi} \cdot (1.136 - 2.53 \chi + 5.6 \chi^2),$$
окуда момент инерции будет равен:
$$I = \bar{I}  M'  R_{eq}^2,$$
а момент импульса находится как:
$$J = I\Omega_{\star}.$$
Гравитация на экваторе же будет равна:
$$g_0 = \frac{GM'}{R_{eq}^2\sqrt{1 - 2 \chi}}$$

Также нам понадобится Кеплеровская (1-ая космическая) скорость на экваторе, равная:
$$V_{kep} =  \sqrt{g_0 R_{eq}},$$
откуда Кеплеровская циклическая частота будет равна:
$$\Omega_{kep} = \frac{V_{kep}}{R_{eq}}.$$

\section{Поверхность нейтронной звезды}
Поверхность нейтронной звезды характеризуют следующие параметры:
\begin{itemize}
    \item $\phi$ - долгота
    \item $\theta$ - полярный угол
    \item $\Omega$ - циклическая частота вращения 
\end{itemize}
Частота вращения $\Omega$ явлется функцией полярного угла $\theta$, и задается нами. Это и есть наша модель слоя растекания, о чем поговорим далее.
\subsection{Метрические коэффициенты}
Для начала, используя уравнение состояния, находим локальный радиус поверхности как:
$$R = R_{eq} \cdot \left[1 + \sum_{l=0,2,4} a_l \cdot P_l(\cos\theta) \right],$$
где коэффициенты $a_l$ равны:
$$a_0 = -0.18 \cdot \bar{\Omega}_{\star}^2 + 0.23 \cdot \chi \bar{\Omega}_{\star}^2 - 0.05 \cdot \bar{\Omega}_{\star}^4,$$
$$a_2 = -0.39 \cdot \bar{\Omega}_{\star}^2 + 0.29 \cdot \chi \bar{\Omega}_{\star}^2 - 0.13 \cdot \bar{\Omega}_{\star}^4,$$
$$a_4 = +0.04 \cdot \bar{\Omega}_{\star}^2 - 0.15 \cdot \chi \bar{\Omega}_{\star}^2 + 0.07 \cdot \bar{\Omega}_{\star}^4,$$
а $P_l(\cos\theta)$ -- полиномы Лежандра. Пересчитаем метрический коэффициент $\bar{\Omega}$ с учетом переменности угловой скорости $\Omega$:
$$\bar{\Omega} =  \Omega \cdot \sqrt{\frac{R_{eq}^3}{G M'}}$$

Для нахождения коэффициента метрики $\bar{r}$, методом последовательных приближений обрабатываем следующую схему:
$$\bar{r} = \frac{R}{e^{-\nu} B},$$
где $\nu$ и $B$ соответсвенно равны:
$$\nu = \nu_0 + \left(\frac{b_c}{3} - q_c P_2(\cos \theta)\right)\bar{u}^3 ,$$
$$B = B_0 + b_c \cdot \bar{u}^2 ,$$
а $\nu_0$ и $B_0$ в свою очередь:
$$\nu_0 = \ln \left(\frac{1 - \frac{\bar{u}}{2}}{1 + \frac{\bar{u}}{2}} \right) ,$$
$$B_0 = \left(1 - \frac{\bar{u}}{2}\right) \left(1 + \frac{\bar{u}}{2}\right).$$
Коэффициенты $q_c$ и $b_c$ равны:
$$q_c = -0.11 \cdot \left(\frac{\bar{\Omega}}{\chi}\right)^2,$$
$$b_c = 0.4454 \cdot \bar{\Omega}^2  \chi,$$
а величина $\bar{u}$ определяется как:
$$\bar{u} =  \frac{R_s}{2\bar{r}},$$
замыкая тем самым систему. Начать метод последовательных приближений можно с $\bar{r}^{(0)} = R$. 
После достижения нужной точности сходимости результата, расчитываем по тем же формулам $\bar{u}$, $\nu$, $B$, а также $\zeta$, равное:
$$\zeta = \zeta_0 + b_c\left( \frac{4}{3} P_2(\cos \theta) - \frac{1}{3} \right)\bar{u}^3,$$
где $\zeta_0$ равно:
$$\zeta_0 = \ln B_0.$$ 
C помощью этих коэффициентов можно вычислить локальный угловой момент $\varpi$, равный:
$$\varpi = \frac{2GJ}{c^2  \bar{r}^3} \cdot (1 - 3 \bar{u})$$
а также релятивистские факторы:
$$\beta' =  \frac{R \varpi}{c} e^{-\nu}  \sin \theta $$
$$\beta =  \frac{R}{c} e^{-\nu} (\Omega - \varpi) \sin \theta $$
$$\gamma = \frac{1}{1 - \beta^2}$$


\subsection{Локальная гравитация}
Локальное ускорение свободного падения расчитывается как:
$$g = g_0 \left[ 1 + e_e \cdot \sin^2\theta + e_p \cdot \cos^2 \theta + e_{60} \cdot \cos \theta \right],$$
где коэффициенты, зависящие от $\bar{\Omega}_{\star}$ равны:
$$e_e = c_e \bar{\Omega}_{\star}^2 + d_e \bar{\Omega}_{\star}^4 + f_e \bar{\Omega}_{\star}^6$$
$$e_p = c_p \bar{\Omega}_{\star}^2 + (d_p - d_{60}) \bar{\Omega}_{\star}^4 + f_p \bar{\Omega}_{\star}^6$$
$$e_{60} = d_{60} \bar{\Omega}_{\star}^4$$
а коэффициенты, зависящие от $\chi$ равны:
$$c_e = 0.776 \chi - 0.791$$
$$c_p = 1.138 - 1.431 \chi$$ 
$$d_e = \chi (2.431 \chi - 1.315) $$
$$d_p = \chi(0.653 - 2.864 \chi)$$
$$d_{60} = \chi (13.47 - 27.13 \chi)$$
$$fe = -1.172 \chi$$
$$fp = 0.975 \chi$$
Однако, на самом деле ускорение свободного падения будет меньше чем $g$, поскольку центробежная сила, связанная с $\bar{\Omega}$, увеличилась.
Поскольку g представляет собой:
$$g = g_0 + a_c + a_{EoS},$$
где $a_{EoS}$ связано с вращением поверхности нейтронной звезды с угловой скоростью $\Omega_{\star}$, а $a_c$ равно:
$$a_c (\bar{\Omega}) = - g_0 \bar{\Omega}^2 \sin^2 \theta \cdot \lambda(\chi),$$
где $\lambda(\chi)$ равна:
$$\lambda(\chi) = \left[1 + \chi(-1 + 2\bar{I}) + \chi^2 (-2 + 4\bar{I} - 8\bar{I}^2)\right]$$
то $g_{eff}$ можно расчитать как: 
$$g_{eff} = g - g_0 (\bar{\Omega}^2 - \bar{\Omega}_{\star}^2 ) \sin^2 \theta \cdot \lambda(\chi)$$
Локальная Эддингтоновская светимость при такой гравитации будет равна:
$$F^{(Edd)} = \frac{cg_{eff}}{\kappa}$$

\subsection{Следствия вращения}
Угол между локальным радиус-вектором $\hat{r}$ и линией наблюдения $\hat{k}$ равен:
$$\cos \psi = \cos i \cos \theta + \sin i  \sin \theta \cos \phi$$
Рассчитаем коэффициент $G_{yu}$, равный:
$$ G_{yu} = 1 + \frac{u^2y^2}{112} -  \frac{\mathrm{e}}{100} uy \cdot \left[\ln \left(1 - \frac{y}{2}\right) + \frac{y}{2}\right],$$
где $y = 1 - \cos \psi$, $u = R_{s}/R$. Теперь мы можем расчитать угол $\alpha$, как:
$$\cos\alpha = 1 - y(1 - u) G_{yu}$$
Теперь рассчитаем фактор линзирования:
$$D = 1 + \frac{3 u^2 y^2}{112} - \frac{\mathrm{e}}{100} uy \left[ 2 \ln\left(1 - \frac{y}{2}\right) + y\frac{1 - 3y/4}{1 - y/2}\right]$$
Найдем ещё несколько углов:
$$\cos \mu = \frac{\cos i - \cos \theta \cos \psi}{\sin \theta \sin \psi} $$

$$\cos \xi = - \frac{\sin \alpha}{\sin \psi} \sin i  \sin\phi $$
$$\cos \sigma = \cos \eta \cos \alpha + \sin \eta \sin \alpha \cos \mu $$
где угол $\eta$ находится как:
$$\cos \eta =  \frac{1}{\sqrt{1 + f^2}},$$
где фактор $f$ равен:
$$f = \frac{1}{\sqrt{1 + R'_s/R}} \frac{1}{R} \frac{dR}{d\theta}$$
а $R'_s = 2GM'/c^2$ -- скорректированный радиус Шврацшильда. Производную радиуса же легко найти аналитически как:

$$\frac{dR}{d\theta} = R_{eq} \cdot \left[1 + \sum_{l=2,4} a_l \cdot \frac{d P_l(\cos\theta)}{d\theta} \right],$$

Допплеровский фактор $\delta$ будет равен:
$$\delta = \frac{1}{\gamma (1 - \beta \cos \xi)}$$
Откуда корректная $\sigma'$ будет равна:
$$\cos \sigma' = \delta \cos \sigma $$.


Элемент площади поверхности будет равен:
$$dS = \frac{1}{\cos \eta} R^2 d\cos\theta d\phi$$

Наблюдаемый телесный угол будет равен:
$$d\Omega_{obs} =  \frac{dS}{d^2} \cos \sigma D $$ 

\section{Спектр нейтронной звезды}

Спектр нейтронной звезды характеризуется парметрами:
\begin{itemize}
    \item $E$ - энергия
    \item $f$ - локальный относительный поток
    \item $T_{eff}$ - локальная эффективная температура
\end{itemize}

Реальная энергия $E'$, на которой мы видим наш спектр расчитывается как:
$$E' = \frac{E}{\delta e^{\nu}(1 + \beta' \cos \xi)}$$
Спектральная плотность энергии $\rho$ явлется результатом интерполяции теоретических спектров на нашу сетку:
$$\rho = \rho_{int}(f, \log(g_{eff}), E')$$
Интенсивность излучения расчитывается как:
$$I = \rho \cdot (0.4215 + 0.86775 \cos \sigma')$$
Тогда наблюдаемый спектр равен:
$$B = \int_{\cos \sigma > 0} \left(\frac{E}{E'}\right)^3 I d\Omega_{obs}$$

Этот спектр в свою очередь может быть приближен разбавленным черным телом и представлен в форме:
$$B \approx w \pi B_{dbb} (f_c T_{eff})$$



\end{document}